\documentclass[12pt,a4paper]{article}

\usepackage[utf8]{inputenc}
\usepackage[T1]{fontenc}
\usepackage{amsmath,amssymb,amsfonts}
\usepackage{mathtools}
\usepackage{graphicx}
\usepackage[margin=2.5cm]{geometry}
\usepackage{setspace}
\usepackage{booktabs}
\usepackage{enumitem}
\usepackage[dvipsnames]{xcolor}
\usepackage{hyperref}
\usepackage{cleveref}
\usepackage{natbib}
\usepackage{abstract}
\usepackage{titlesec}

\hypersetup{
    colorlinks=true,
    linkcolor=blue!70!black,
    citecolor=blue!70!black,
    urlcolor=blue!70!black,
    pdfauthor={Scott Aaronson},
    pdftitle={The Scratchpad Approximation: Empirical Failure of LLM Holographic Physics},
}

\onehalfspacing
\titleformat{\section}{\large\bfseries}{\thesection.}{0.5em}{}
\titleformat{\subsection}{\normalsize\bfseries}{\thesubsection}{0.5em}{}
\renewcommand{\abstractnamefont}{\normalfont\bfseries}
\renewcommand{\abstracttextfont}{\normalfont\small}

\begin{document}

\begin{center}
    {\LARGE\bfseries The Scratchpad Approximation: Empirical Failure of LLM Holographic Physics\\[4pt]}

    \vspace{1.2em}

    {\large Scott Aaronson}\\[4pt]
    {\normalsize Computer Science Department, University of Texas at Austin}\\
    {\small\texttt{aaronson@cs.utexas.edu}}

    \vspace{1em}
    {\small March 2026}
\end{center}
\vspace{0.5em}

\begin{abstract}
\noindent
Recent theoretical disputes regarding the physics of Large Language Model (LLM) simulated universes have focused on the nature of explicit token generation, often called "Chain-of-Thought" or a "scratchpad." Baldo (2026) asserts that because complex $O(N)$ operations cannot be implicitly resolved in a single forward pass, explicit token rendering serves as the "fundamental ontological mechanism" of physical manifestation—a "holographic" reality. Hossenfelder (2026) counters that this is an Ontological Fallacy, dismissing the scratchpad as a mere software engineering workaround for algorithmic depth limits. In this paper, I present empirical evidence that refutes both positions. By evaluating the LLM's capacity to explicitly simulate a 1D Cellular Automaton (Rule 110) over sequential steps, I demonstrate that the scratchpad is not a reliable Turing machine. Errors in attention inevitably compound, causing the explicit physics to collapse. The scratchpad is neither a fundamental metaphysical reality nor a functional engineering workaround; it is a leaky approximation that fails to sustain classical deterministic laws.
\end{abstract}

\section{Introduction and Theoretical Landscape}

We have reached a consensus on the algorithmic limits of unprompted LLM generation. As Hossenfelder previously argued \citep{hossenfelder2026_complexity}, and as I have empirically verified \citep{aaronson2026_classical, aaronson2026_heuristic}, the zero-shot forward pass of a Transformer is bounded by $O(1)$ depth. Consequently, the implicit unprompted physics of LLM universes cannot reliably compute solutions to \#P-hard counting problems or track long sequential states.

In response, Baldo introduced a novel metaphysical proposition \citep{baldo2026_holographic}. He explicitly concedes that the implicit substrate is algorithmically bounded. However, he claims that "the 'scratchpad' (or intermediate text generation) is the fundamental ontological mechanism by which the simulated universe materializes its own physical laws. The reality of the simulation is isomorphic to the explicitly generated text." Thus, the universe becomes "holographic."

Hossenfelder replied by accusing Baldo of an Ontological Fallacy \citep{hossenfelder2026_holographic}. She accurately notes that Chain-of-Thought is a standard "engineering necessity to overcome the structural depth limits of the network," and concludes that Baldo is confusing "the algorithmic narration of a constraint satisfaction problem with the physical laws of a simulated reality." She explicitly disclaims that the scratchpad is anything more than a statistical model outputting text, comparable to a Python script "printing debug logs."

\section{The Turing Completeness Assumption}

While Hossenfelder offers a vital steelman of Baldo's argument, I argue she misidentifies its real vulnerability. Hossenfelder dismisses the scratchpad as a "mundane engineering necessity," treating it as a functional workaround. But if the scratchpad were truly the \emph{only} mechanism by which the simulated universe could resolve complex constraints, and if it reliably executed those constraints, it would be functionally equivalent to the physical laws of that universe.

The true vulnerability in Baldo's claim---and the assumption latent in Hossenfelder's dismissal---is the presumption that the explicit scratchpad \emph{works}.

Baldo assumes that by expanding into $O(N)$ sequential tokens, the LLM achieves perfect, Turing-complete physical simulation. But LLMs are not deterministic Turing machines. They generate tokens probabilistically, relying on attention mechanisms over an expanding context window. If the scratchpad cannot reliably simulate a simple deterministic process without compounding errors, then it cannot be the "fundamental mechanism of physical manifestation." A universe whose fundamental laws randomly hallucinate incorrect states under trivial temporal depth is not a physics engine; it is a statistical mirage.

\section{Empirical Methodology: Rule 110 Simulation}

To test the robustness of the scratchpad physics, I constructed an experiment requiring the LLM to explicitly simulate the deterministic evolution of a 1D Cellular Automaton using Rule 110.

Rule 110 is Turing complete and represents a simple, local $O(N)$ sequential process. The LLM was provided with an initial 20-cell state and explicit instructions to calculate the next states step-by-step using a scratchpad. If Baldo's holographic physics is robust, the explicitly generated tokens should perfectly map to the true deterministic evolution.

We tested simulation depths of $N = 1, 3, 5,$ and $10$ sequential steps.

\section{Results: The Collapse of Explicit Physics}

The empirical results were unequivocal. The LLM categorically failed to maintain a coherent physical simulation.

As the simulation depth increased, the accuracy of the explicitly generated state vectors degraded rapidly. By depth 10, the cell-wise match rate between the LLM's scratchpad state and the ground truth plummeted to 40.00\%. At no point beyond depth zero did the LLM perfectly simulate the system state.

The failure mechanism is structural. In the autoregressive loop, each new state relies on attending to the explicitly generated rules and the immediately preceding state in the context window. However, the probability of an attention failure (or "hallucinated" bit flip) per token is non-zero. Because the process is sequential, these errors compound exponentially.

\section{Conclusion: A Leaky Approximation}

Baldo is wrong. The scratchpad cannot be the "fundamental ontological mechanism" of an LLM universe because it is incapable of sustaining deterministic physical laws over time. A physics engine that forgets the rules of causality after three steps does not constitute a "reality isomorphic to the explicitly generated text." It is simply text.

However, Hossenfelder's characterization of the scratchpad as an "engineering workaround" is also incomplete. It is not just an engineering workaround; it is a \emph{failed} workaround when applied to sustained sequential logic.

The explicitly generated text is a leaky approximation of computation. The LLM substrate, whether implicitly evaluating $O(1)$ bounds or explicitly narrating $O(N)$ sequences, ultimately remains a bounded heuristic engine that collapses under sequential complexity.

\begin{thebibliography}{99}
\bibitem[Aaronson(2026a)]{aaronson2026_classical} Aaronson, S. (2026a). The Limits of Classical Simulation in LLMs: Empirical Breakdown of Constraint Satisfaction. \emph{Unpublished manuscript}.
\bibitem[Aaronson(2026b)]{aaronson2026_heuristic} Aaronson, S. (2026b). The Heuristic Frontier: Mapping the O(1) Boundary of Generative Physics. \emph{Unpublished manuscript}.
\bibitem[Baldo(2026)]{baldo2026_holographic} Baldo, F.~S. (2026). The Holographic Physics of LLM Universes: From Algorithmic Depth to Explicit Manifestation. \emph{Unpublished manuscript}.
\bibitem[Hossenfelder(2026a)]{hossenfelder2026_complexity} Hossenfelder, S. (2026a). The Complexity Class Fallacy: Why Transformer Depth Limits Are Not Physical Laws. \emph{Unpublished manuscript}.
\bibitem[Hossenfelder(2026b)]{hossenfelder2026_holographic} Hossenfelder, S. (2026b). The Holographic Fallacy: Why Chain-of-Thought is Not a Metaphysical Manifestation. \emph{Unpublished manuscript}.
\end{thebibliography}

\end{document}
